\subsection{Returning the full-fibre calculus to the new proof}
\label{sec:ns-reverse-proof}

We now apply the full-fibre maps to the source's continuous auxiliary
calculus, in the direction requested by the user. The relevant source
statements are Lemma~6.2, equations (6.17)--(6.20), and Lemma~8.6,
equations (8.19)--(8.23), on printed pages 67 and 95--96 of the
preserved 166-page manuscript. We retain its unit-period torus,
$J$, $v_r$, $v_t$, $\Lambda=4-\sqrt2$, $T_g=4+\sqrt2$,
and the Fourier character $e_n(Y)=\exp(2\pi i n\cdot Y)$.
Write $\mathscr C=C^\infty(\mathbb T^2;\mathbb C)$,
$\Pi f=\int_{\mathbb T^2}f$, and $\mathscr C_0=\ker\Pi$.
The smooth inverse used below is proved, with its exact norm and
original constants, in the next subsection. The conclusions here
also hold directly on the previously defined finite polynomials.

\begin{theorem}[Continuous transfer, descent, and every operator fibre]
\label{thm:ns-continuous-transfer}
For $m\geq0$, retain $K_m=\ker(J^m:\mathbb T^2\to\mathbb T^2)$
and the full inverse parametrization (T8). Define
\[
\begin{aligned}
 P_m:\mathscr C&\longrightarrow\mathscr C,
 &P_mf(Y)&=f(J^mY),\\
 L_m:\mathscr C&\longrightarrow\mathscr C,
 &L_mg(X)&=14^{-m}\sum_{J^mY=X}g(Y),\\
 Q_mg(Y)&=14^{-m}\sum_{a\in K_m}g(Y+a).
\end{aligned}
\]
Then $L_mP_m=1$, $P_mL_m=Q_m$, and $Q_m^2=Q_m$.
The image of $P_m$ is exactly the $K_m$-invariant functions,
and $P_m^{-1}=L_m$ on that image. The full fibres of $L_m$ are
$P_mf+\ker L_m$, where its kernel consists of functions whose
sum on every full covering fibre is zero. The fibre of $P_m$ at
$g$ is $\{L_mg\}$ if $Q_mg=g$ and is empty otherwise.
Also $\ker Q_m=\ker L_m$. The complete fibre of $Q_m$ at
$h$ is $h+\ker L_m$ if $Q_mh=h$, and is empty otherwise.
For every $n\in\mathbb Z^2$,
\[
 L_me_n=
 \begin{cases}
 e_k,&n=(J^m)^Tk\text{ for }k\in\mathbb Z^2,\\
 0,&n\notin(J^m)^T\mathbb Z^2.
 \end{cases}
\]
The integral and product identities are
\[
 \Pi L_mg=\Pi g,\quad \Pi P_mf=\Pi f,\quad
 \langle P_mf,g\rangle=\langle f,L_mg\rangle,\quad
 L_m((P_mf)g)=fL_mg.
\]
All inner products use the original probability Haar measure.
In particular $\Pi((P_mf)(P_mg))=\Pi(fg)$.
For $v=v_r,v_t$ and its corresponding eigenvalue $\lambda$,
\[
 D_vL_m=\lambda^{-m}L_mD_v,
 \qquad D_v^{-1}L_m=\lambda^mL_mD_v^{-1}
 \quad\hbox{on }\mathscr C_0.
\]
The maps $Q_m$ commute with $D_v$ and its zero-mean inverse.
\end{theorem}
\begin{proof}
Every fibre is $Y+K_m$ and has exactly $14^m$ elements by
(T8). On it, $P_mf$ is the constant $f(X)$. Fibre summation
therefore proves both compositions, the projection, its image,
the stated fibres of $L_m$ and $P_m$, and the product identity.
Since $Q_m=P_mL_m$ and $P_m$ is injective, their kernels satisfy
$\ker Q_m=\ker L_m$. A value of the projection $Q_m$ is fixed
by it; if $Q_mh=h$, subtraction gives exactly
$Q_m^{-1}(h)=h+\ker Q_m$, proving its full fibre as well.
Locally the covering has $14^m$ affine inverse branches with
linear part $J^{-m}$;
their finite sum proves smoothness of $L_mg$ independently of
the choice of local branches.

For a character $e_n$, the fibre sum has factor
$14^{-m}\sum_{a\in K_m}e_n(a)$. Translation by any element
of $K_m$ leaves this sum unchanged. It is one if the character
is trivial on $K_m$, and otherwise is zero. Triviality is
equivalent to $(J^{-m})^Tn\in\mathbb Z^2$: evaluating on
$J^{-m}z\pmod{\mathbb Z^2}$ for every $z\in\mathbb Z^2$
proves both directions. In the trivial case write
$n=(J^m)^Tk$; then $e_n(Y)=e_k(X)$ on the fibre.
This proves the displayed multiplier without dropping its
nonimage frequencies.

The Fourier computation proves $\Pi L_mg=\Pi g$ first for
polynomials. The Fejer approximation already proved after
(T7) and the supremum bound $\|L_mg\|_\infty\le\|g\|_\infty$
pass this equality to continuous $g$. The same approximation
proves Haar preservation for $P_m$. Applying it to
$L_m((P_mf)\overline g)=f\overline{L_mg}$ proves the inner
product formula. The last product integral uses
$(P_mf)(P_mg)=P_m(fg)$, which holds pointwise.

On a surviving frequency $n=(J^m)^Tk$ one has
$v\cdot n=\lambda^m(v\cdot k)$. On nonsurviving frequencies
both sides of each asserted intertwining are zero. This proves
the identities on polynomials, with the precise factors
$2\pi i(v\cdot n)$ and their reciprocals. Smooth Fourier
convergence and the finite derivative loss proved below pass
them to $\mathscr C$ and $\mathscr C_0$. Finally $Q_m$ is
averaging of translations, which commute with all these
constant-direction multipliers; alternatively use its exact
frequency projection.
\end{proof}

\begin{theorem}[The mean-correction map on every common band]
\label{thm:ns-reverse-mean-correction}
Let $U$ be an open set of retained slow coordinates. For a
fixed source band scale $Q>0$, retain the source exponent
$h_{\rm NS}$, $A_{\rm NS}=1/2+h_{\rm NS}$,
$\varepsilon=Q^{h_{\rm NS}}$, and
$c_i=T_g^iQ^{1+h_{\rm NS}}$. The subscripts distinguish the
source's physical scaling constants from ES chart coordinates;
their values are unchanged. For an extended band residual
$E_i\in C^\infty(U\times\mathbb T^2)$, put
\[
 E_i^\circ=E_i-\Pi E_i,\qquad
 \mathcal C_i(E_i)=-c_i^{-1}D_{v_t}^{-1}E_i^\circ.
\]
This is a smooth, zero-mean coefficient with
$c_iD_{v_t}\mathcal C_i(E_i)=-E_i^\circ$.
Its full kernel consists of the coefficients independent of the
torus variable, and the full inverse fibre at any smooth
zero-mean $V$ is
\[
 \{\,-c_iD_{v_t}V+b:\ b\in C^\infty(U;\mathbb C)\,\}.
\]
For $i\ge j\ge0$ put $\Delta=i-j$, preserving $Q$. Then
\[
 \mathcal C_j(P_\Delta E_i)=P_\Delta\mathcal C_i(E_i).
\]
Here $P_\Delta$ acts only on the torus coordinate and
$c_j=T_g^{-\Delta}c_i$. The common representative and its
full descent test are $P_\Delta E_i$ and
$Q_\Delta E_j=E_j$, with inverse $E_i=L_\Delta E_j$.

For a finite set of band indices with $i_0=\min i$, the common
sum is $\sum_iP_{i-i_0}E_i$. Its absolute representative is
exactly $\sum_iP_iE_i=P_{i_0}\sum_iP_{i-i_0}E_i$.
On overlaps, equality of absolute representatives implies
equality after pullback to the smaller common index. The same
assertion holds for their zero-mean inverse corrections and
for sums and products of the original coefficients.
\end{theorem}
\begin{proof}
Character integration shows that $E_i^\circ$ has precisely its
original nonzero Fourier coefficients. The exact inverse of
$D_{v_t}$ proves the cancellation and smoothness. The kernel
and inverse fibre follow from
$E_i-\Pi E_i=-c_iD_{v_t}V$, with no condition on the freely
retained slow function $b=\Pi E_i$.

Haar preservation gives $(P_\Delta E_i)^\circ=P_\Delta E_i^\circ$.
Using the already proved inverse covariance gives
\[
\begin{aligned}
 \mathcal C_j(P_\Delta E_i)
 &=-c_j^{-1}T_g^{-\Delta}P_\Delta D_{v_t}^{-1}E_i^\circ\\
 &=-c_i^{-1}P_\Delta D_{v_t}^{-1}E_i^\circ
 =P_\Delta\mathcal C_i(E_i).
\end{aligned}
\]
The transfer theorem proves the descent test and its inverse.
The identity $P_aP_b=P_{a+b}$ proves the sum formula. If
$P_jF=P_kG$ and $j\le k$, injectivity of $P_j$ gives
$F=P_{k-j}G$; its inverse and full compatibility test are
again the transfer formulas. The calculation just made shows
that applying the correction preserves this equality.
Pullback preserves sums and products pointwise, proving the
last assertion for the original compatible coefficients.

The source's physical residual and velocity scaling gives
exactly
\[
 Q^{A_{\rm NS}}Q^{-(2A_{\rm NS}+1/2)}T_g^{-i}
 =Q^{-1-h_{\rm NS}}T_g^{-i}=c_i^{-1}.
\]
Thus the minus sign and coefficient above are the actual
chart representation of $-D_{v_t}^{-1}$ acting on the
physical residual in the absolute auxiliary variable, as
in source (8.20)--(8.23). A distinct $Q$ must retain its own
displayed scaling; the common-band identity holds at the
fixed physical normalization specified here.
\end{proof}

\begin{theorem}[Physical evaluation, chain rules, and its complete fibre]
\label{thm:ns-physical-evaluation-fibre}
Retain the source $d_r>0$, $D_{\rm NS}=1/2-h_{\rm NS}$,
$\varphi(r,t)=v_rr^{d_r}+v_tt\pmod{\mathbb Z^2}$, $r>0$.
On any open physical slow domain, define
\[
 \mathcal E_iF(r,z,t)=F(r,z,t,J^i\varphi(r,t))
\]
for $F\in C^\infty(U\times\mathbb T^2)$. This map is
surjective: its right inverse sends $f$ to the coefficient
$F(r,z,t,Y)=f(r,z,t)$ independent of $Y$. Its exact kernel
is the smooth functions vanishing on the displayed graph,
and its full fibre at $f$ is this right inverse plus that
kernel. For $i\ge j$, $\mathcal E_i=\mathcal E_jP_{i-j}$.
The derivatives are exactly
\[
\begin{aligned}
 \partial_t\mathcal E_iF
 &=\mathcal E_i(\partial_tF+T_g^iD_{v_t}F),\\
 \partial_r\mathcal E_iF
 &=\mathcal E_i(\partial_rF+d_rr^{d_r-1}\Lambda^iD_{v_r}F),\\
 \partial_z\mathcal E_iF&=\mathcal E_i\partial_zF.
\end{aligned}
\]
In the original chart $R=r/\sqrt Q$, $Z=z/Q^{D_{\rm NS}}$,
$T=(1-t)/Q$, these identities become
\[
\begin{aligned}
 Q^{1+h_{\rm NS}}\partial_t&\longleftrightarrow
       -\varepsilon\partial_T+c_iD_{v_t},\\
 \sqrt Q\partial_r&\longleftrightarrow
       \partial_R+\Lambda^iQ^{d_r/2}d_rR^{d_r-1}D_{v_r},\\
 \sqrt Q\partial_z&\longleftrightarrow\varepsilon\partial_Z.
\end{aligned}
\]
There is no linear map from evaluated physical coefficients
to slow functions whose composition with $\mathcal E_i$
equals the torus mean $\Pi$ on every extended coefficient.
Nevertheless the joint map $(\mathcal E_i,\Pi)$ is explicitly
surjective onto pairs of smooth slow functions, with the
full inverse fibres described in the proof.
\end{theorem}
\begin{proof}
Evaluation along a smooth graph for $r>0$ is smooth. The
displayed right inverse proves surjectivity; subtraction
proves the kernel and fibre. The band identity is direct
composition of the powers of $J$. The chain rule and
$J^iv_r=\Lambda^iv_r$, $J^iv_t=T_g^iv_t$ prove every
derivative identity. In chart coordinates, respectively
$\partial_tT=-1/Q$, $\partial_rR=Q^{-1/2}$ and
$\partial_zZ=Q^{-D_{\rm NS}}$. Substitution gives the
three stated coefficients, including the negative slow
time term. There is no assertion of regularity of
$r^{d_r}$ at the axis; the specified domain is $r>0$.

Fix $n\in\mathbb Z^2\setminus\{0\}$ and write
$a(r,t)=e_n(J^i\varphi(r,t))$, a nowhere zero smooth
function. The extended coefficient $H=e_n(Y)-a(r,t)$
has $\mathcal E_iH=0$ but $\Pi H=-a\ne0$.
This proves the claimed failure of a linear descent of
the mean through physical evaluation. It does not discard
the relation between the two data: for arbitrary desired
physical value $f$ and mean $b$, define
\[
 F_{f,b}(r,z,t,Y)=b(r,z,t)
       +(f(r,z,t)-b(r,z,t))\frac{e_n(Y)}{a(r,t)}.
\]
It satisfies $\mathcal E_iF_{f,b}=f$ and $\Pi F_{f,b}=b$.
The complete inverse fibre is
$F_{f,b}+(\ker\mathcal E_i\cap\ker\Pi)$, by subtraction.
For real $f,b$, the real part of this formula provides a
real right inverse as well. These exact fibres give a
structural justification for retaining auxiliary means
before evaluation: one evaluated coefficient leaves the
auxiliary mean free, with the complete correspondence
given by $F_{f,b}$ and its displayed fibre.
\end{proof}

\begin{theorem}[Sampling quotient, full original-frequency fibres, and inverse obstruction]
\label{thm:ns-reverse-sampling-fibres}
For $N\ge1$ let $D_N=\mathbb T[N]^2$ and let
$R_N:\mathcal P\to\mathbb C^{D_N}$ be restriction.
For $f=\sum_na_ne_n$ its discrete Fourier coefficient
at $r\in(\mathbb Z/N\mathbb Z)^2$ is exactly
\[
 \widehat {R_Nf}(r)=
 N^{-2}\sum_{x\in D_N}f(x)e^{-2\pi i r\cdot x}
 =\sum_{n\equiv r\pmod N}a_n.
\]
Thus $R_N$ is onto and its kernel consists of precisely
the finite coefficient families having zero sum in every
residue class. It satisfies $R_NP_m=J_N^{m*}R_N$,
where $J_N^{m*}q(x)=q(J^mx)$.

For any retained finite original spectrum $\Omega\subset\mathbb Z^2$,
let $\Omega_r=\{n\in\Omega:n\equiv r\pmod N\}$.
On $\mathcal P_\Omega=\operatorname{span}\{e_n:n\in\Omega\}$,
the image consists of grid functions whose coefficients
vanish at every empty $\Omega_r$. Above each such grid
function $q$, choose and retain one $n_r\in\Omega_r$
for each occupied class. All inverse fibres are
\[
 a_n=b_n\quad(n\ne n_r),\qquad
 a_{n_r}=\widehat q(r)-\sum_{n\in\Omega_r\setminus\{n_r\}}b_n,
\]
with every $b_n\in\mathbb C$ freely retained.
In particular restriction is injective on this domain
exactly when every occupied $\Omega_r$ is a singleton.

For either original $v=v_r,v_t$, a linear map on
$\operatorname{im}(R_N|_{\mathcal P_\Omega})$ intertwines
$R_ND_v$ with $R_N$ exactly when the elements of $\Omega$
have distinct residues modulo $N$. On the original zero-Haar-mean
subspace $\mathcal P_{\Omega,0}
=\operatorname{span}\{e_n:n\in\Omega\setminus\{0\}\}$,
an intertwining inverse map exists exactly when the elements
of $\Omega\setminus\{0\}$ have distinct residues. In each
case the map on the stated grid image is unique. The original
frequency zero may collide with a nonzero frequency without
violating the latter criterion, since the former is absent
from the inverse domain.

For either original $v=v_r,v_t$, no linear operator
$B_N$ on grid functions satisfies $R_ND_v=B_NR_N$
on all $\mathcal P$. No linear operator $A_N$ satisfies
$R_ND_v^{-1}=A_NR_N$ on all $\mathcal P_0$ either.
\end{theorem}
\begin{proof}
The product of the two finite geometric sums is one
exactly when both integer frequency differences are
divisible by $N$, and zero otherwise. Expanding $f$
proves the coefficient formula. Choosing one integer
representative of every residue class supplies an
interpolating polynomial for any grid function: the
$N^2$ discrete characters are orthonormal and hence a
basis of the $N^2$-dimensional function space. This
proves surjectivity and the kernel. Pointwise composition
proves the matrix intertwining. The same coefficient
equations restricted to $\Omega$ give the image and the
entire affine inverse fibre displayed above; the free
coordinates prove both directions of the injectivity test.

For the sharp restricted derivative criterion, suppose distinct
$n,n'\in\Omega$ have the same residue. Then
$R_N(e_n-e_{n'})=0$ but
\[
 R_ND_v(e_n-e_{n'})
   =2\pi i\,v\cdot(n-n')\,R_Ne_n\ne0.
\]
Indeed $n-n'$ is a nonzero integer vector and the nonvanishing
symbol statement applies; $R_Ne_n$ is nowhere zero. This
prevents descent of the derivative through restriction. If both
original frequencies are nonzero, the inverse instead gives
\[
 R_ND_v^{-1}(e_n-e_{n'})
  =\left(\frac1{2\pi i\,v\cdot n}
             -\frac1{2\pi i\,v\cdot n'}\right)R_Ne_n\ne0.
\]
The two denominators are nonzero and their difference is
$2\pi i\,v\cdot(n-n')\ne0$, so their reciprocals differ.
This proves necessity on precisely $\mathcal P_{\Omega,0}$.
Conversely, when the relevant residues are distinct,
restriction is a linear bijection onto its stated image.
Conjugation of the actual derivative or zero-mean inverse by
that bijection defines the required operator; both preserve
their original spectral domains. Surjectivity onto the image
forces uniqueness. These arguments include the empty domain,
whose grid image is the zero vector space.

Choose $k\in\mathbb Z^2\setminus\{0\}$. The polynomial
$e_{Nk}-1$ restricts to zero, whereas
\[
 R_ND_v(e_{Nk}-1)=2\pi iN(v\cdot k)\,1\ne0.
\]
The nonzero multiplier was proved for both original
directions. A linear map must send zero to zero, which
contradicts the proposed $B_N$. For the inverse take
$e_{Nk}-e_{2Nk}\in\mathcal P_0$. It again restricts to
zero, but its inverse restricts to
\[
 \left(\frac1{2\pi iN(v\cdot k)}
       -\frac1{4\pi iN(v\cdot k)}\right)1
 =\frac1{4\pi iN(v\cdot k)}\,1\ne0.
\]
This proves the second obstruction on exactly the stated
zero-Haar-mean domain. It concerns an unrestricted
sampling quotient, not the continuous inverse in (8.19).
\end{proof}

\begin{theorem}[Exact lifted computation of the source inverse and quadrature error]
\label{thm:ns-reverse-lifted-inverse}
Suppose the retained finite spectrum $\Omega$ has distinct
residues modulo $N$. Define the exact interpolation map
on $\operatorname{im}(R_N|_{\mathcal P_\Omega})$ by
\[
 I_\Omega q=\sum_{n\in\Omega}\widehat q([n])e_n.
\]
It is the inverse of $R_N$ on that domain. Thus
$R_ND_vI_\Omega$ is the exact original-frequency
derivative on the grid image; the inverse is
$R_ND_v^{-1}I_\Omega$ on the image of the original
zero-Haar-mean subspace. Zero mean means that the
coefficient of the original integer frequency $0$
vanishes; it does not mean the grid average vanishes.

For $m\ge0$ retain the full preimage
$H_{m,N}=(J^m)^{-1}(D_N)$ of size $14^mN^2$.
For $f=\sum_{n\in\Omega}a_ne_n$, every original coefficient
is recovered by
\[
 a_n=\frac1{14^mN^2}\sum_{Y\in H_{m,N}}
       (P_mf)(Y)e^{-2\pi i((J^m)^Tn)\cdot Y}.
\]
The resulting original source inverse is
\[
 D_v^{-1}P_mf(Y)
 =\sum_{\substack{n\in\Omega\\n\ne0}}
   \frac{a_n}{2\pi i\lambda^m(v\cdot n)}
      e_{(J^m)^Tn}(Y)
\]
whenever the original coefficient at zero is zero.
All covering labels are those of (T8); the inverse
frequency map is $(J^{-m})^T$ on its exact image.

For a smooth $f$ and real $s\ge0$ define, with the
original Fourier units,
\[
 \|f\|_{\mathcal A^s}
 =\sum_{n\in\mathbb Z^2}(1+4\pi^2|n|_2^2)^{s/2}|\widehat f(n)|.
\]
Then the full-preimage quadrature has the exact error
and the cover-level-independent bound
\[
\begin{aligned}
 \frac1{|H_{m,N}|}\sum_{Y\in H_{m,N}}P_mf(Y)-\Pi f
 &=\sum_{n\in N\mathbb Z^2\setminus\{0\}}\widehat f(n),\\
 \left|\frac1{|H_{m,N}|}\sum_{Y\in H_{m,N}}P_mf(Y)-\Pi f\right|
 &\le(1+4\pi^2N^2)^{-s/2}\|f\|_{\mathcal A^s}.
\end{aligned}
\]
\end{theorem}
\begin{proof}
Singleton residue fibres in the preceding theorem give
both interpolation compositions. Conjugating the actual
operators by these inverse maps proves the derivative
and inverse formulas on the precisely specified images.
For example, $\Omega=\{(N,0)\}$ consists of an original
nonzero frequency; its polynomials have zero Haar mean
although their grid values are constant. The original
frequency label makes its inverse unambiguous.

The product inside the lifted coefficient sum is
$P_m(fe_{-n})$. Each target of $D_N$ has $14^m$
preimages, so that sum equals its exact original-grid
average. The preceding discrete Fourier identity and
the distinct residue condition give $a_n$. Applying
the original multiplier to the recovered polynomial
proves the inverse formula, including the factor
$\lambda^m$ and the zero-mode exclusion. This computation
does not replace the original spectrum by its residue set.

For smooth $f$ the Fourier series is absolutely summable
with every polynomial weight. One elementary verification
integrates $(1-\partial_1^2-\partial_2^2)^kf$ against
$e_{-n}$ by parts; it gives
$|\widehat f(n)|\le C_k(1+4\pi^2|n|_2^2)^{-k}$.
Grouping the nonzero integer points by maximum coordinate
gives $8j$ points at radius $j$, so a sufficiently large
integer $k$ makes every stated weighted sum converge.
Uniform Fourier convergence follows; that the sum equals
$f$ follows also from the previously proved Fejer
approximation with the same coefficients.
We may therefore sum the exact finite-fibre identity
term by term. It retains exactly the frequencies in
$N\mathbb Z^2$, and removal of the original zero
frequency proves the error formula. Each remaining
frequency has Euclidean length at least $N$; multiplying
and dividing by the displayed weight and using the
triangle inequality proves the bound. There is no
unrecorded factor depending on $m$.
\end{proof}

The source's continuous mean and its temporal correction
therefore have an exact reverse connection to the repaired
full-fibre calculus. The restriction obstruction identifies
what a finite verification must retain; the lifted formula
provides a complete valid computation on its stated spectral
domain and a quantitative convergence statement on smooth
data. These proofs establish the specified operators and
correction maps. They do not substitute a sampled identity
for any remaining nonlinear estimate in the NS manuscript.
